\documentclass[12pt]{article}
\usepackage[utf8]{inputenc}
\usepackage[T1]{fontenc}
\usepackage{hyperref}
\title{Guinean Sign Language Language Resources for K-12 Curriculum Development}
\author{Compiled by Claude Sonnet 5.5 (claude-sonnet-5-5) for Andr\'as Kornai}
\date{}
\begin{document}
\maketitle

\section{Introduction}
Guinean Sign Language (ISO 639-3 gus) is the sign language of the Deaf community of Guinea (Conakry). Glottolog rates it not endangered and notes a claim, without supporting evidence, that it is related to American Sign Language. Do not confuse it with Língua Gestual Guineense of Guinea-Bissau, which several search results concern. Searches found no grammar, lexicon or school materials for Guinean Sign Language itself.

\section{Primary Educational Resources}
None found.

\section{Linguistic and Reference Materials}
Glottolog language page: \url{https://glottolog.org/resource/languoid/id/guin1250}.

\section{Religious and Community Resources}
Joshua Project profile of the Deaf in Guinea (population figure about 64,000, needs for interpreters in universities and health care): \url{https://www.joshuaproject.net/index.php/people_groups/19007/GV}.

\section{Digital and Audio Resources}
None found.

\section{Curriculum Development Context}
No evidence of a sign-language-medium curriculum was found. Development would require partnership with Guinean Deaf associations and the creation of a documented lexicon and video corpus from scratch.

\end{document}